\documentclass[11pt]{article}

% ============================================================================
%  Audit report: background-subtracted Z4c implementation in AthenaK (src/z4c)
%  Branch audited: project/bg_subtract
% ============================================================================

\usepackage[margin=1in]{geometry}
\usepackage{amsmath,amssymb}
\usepackage{bm}
\usepackage{xcolor}
\usepackage{listings}
\usepackage{booktabs}
\usepackage{longtable}
\usepackage{enumitem}
\usepackage[hidelinks]{hyperref}
\usepackage{titlesec}

\titleformat{\section}{\large\bfseries}{\thesection}{0.6em}{}
\titleformat{\subsection}{\normalsize\bfseries}{\thesubsection}{0.6em}{}

\definecolor{codebg}{gray}{0.96}
\lstset{
  basicstyle=\ttfamily\footnotesize,
  backgroundcolor=\color{codebg},
  breaklines=true,
  columns=fullflexible,
  frame=single,
  framerule=0pt,
  xleftmargin=4pt,
  xrightmargin=4pt,
  language=C++,
  keywordstyle=\color{blue!60!black},
  commentstyle=\color{green!40!black},
  showstringspaces=false
}

\newcommand{\code}[1]{\texttt{#1}}
\newcommand{\file}[1]{\texttt{#1}}
\newcommand{\Lie}{\mathcal{L}}
\newcommand{\dtb}{(\partial_t-\Lie_\beta)}
\newcommand{\gt}{\tilde\gamma}
\newcommand{\At}{\tilde A}
\newcommand{\Gamt}{\tilde\Gamma}
\newcommand{\bg}[1]{#1_{0}}
\newcommand{\res}[1]{#1_{\mathrm s}}

% --- field shorthands ---
\newcommand{\Fop}{\mathcal{F}}            % continuum RHS operator
\newcommand{\Fnum}{\Fop^{h}}              % discrete RHS operator
\newcommand{\half}{\tfrac12}
\newcommand{\third}{\tfrac13}

\title{\bfseries Audit of the Background--Subtracted Z4c Implementation in
AthenaK\\[2pt]
\large \file{src/z4c} on branch \file{project/bg\_subtract}}
\author{Code audit report}
\date{\today}

\begin{document}
\maketitle

\begin{abstract}
This report audits the background--subtracted (residual) Z4c evolution
implemented in \file{src/z4c} of the \file{project/bg\_subtract} branch of
AthenaK, against (i) the full Z4c equations on \file{origin/main}, (ii) the
BSSN background--subtraction method of Lam, Shibata \& Kiuchi
(arXiv:2212.10891; hereafter \textbf{LSK22}), and (iii) the AthenaK Z4c method
paper of Zhu \emph{et al.} (arXiv:2409.10383; hereafter \textbf{Z24}). The
central design property---that the residual right--hand side (RHS) is the
\emph{discrete} difference $\Fnum(\text{full})-\Fnum(\text{background})$, so
that a vanishing residual produces an exactly vanishing RHS to round--off and
the leading truncation error of the (static) background does \emph{not} act as
a source on the residual---is implemented correctly and consistently across
the constraint/evolution sector, the gauge sector, the algebraic--constraint
projection, the Kreiss--Oliger dissipation, the Sommerfeld boundary, and the
AMR prolongation/restriction. The gauge sector is a genuine
\emph{background--adapted} $1+\log$/$\Gamma$--driver that retains the
$1+\log$ principal symbol. We document the term--by--term verification, the
data flow, and a ranked list of residual robustness concerns. We emphasise that
the Kerr--Schild background is an \emph{exact and stable} fixed point of the
residual gauge flow (the moving--puncture damping channels still act on the
residual), so this is not a stability defect; the one genuine trade--off versus
LSK22's $K_0=0$ trumpet is interior regularity---Kerr--Schild is horizon
penetrating and reaches the singularity, so an inner freeze/excision layer is
required where a trumpet would be regular and excision--free.
\end{abstract}

\tableofcontents

% ===========================================================================
\section{Scope, method, and notation}
% ===========================================================================

\paragraph{What was audited.} The diff of \file{src/z4c} between
\file{project/bg\_subtract} and \file{origin/main} (merge base
\code{6ae35f4}). The substantive changes are in
\file{z4c\_calcrhs.cpp} (RHS, $+1118/-845$ effective rewrite),
\file{z4c.cpp} (state containers and background/residual transforms),
\file{z4c.hpp} (data structures and flags),
\file{z4c\_adm.cpp} (ADM$\leftrightarrow$Z4c and constraints on the full
state), \file{z4c\_tasks.cpp} (task ordering), and \file{z4c\_amr.cpp}
(refinement on the reconstructed full state). The driving problem generator is
\file{src/pgen/z4c\_tov\_ks.cpp} (TOV star as a residual on an analytic
Kerr--Schild background). Representative input:
\file{inputs/tde/aurora/z4c\_tov\_ks\_n3\_schwarzschild\_resgauge\_full\_2n\_aurora.athinput}.

\paragraph{Z4c continuum system.} We use the Z4c equations as in Z24 (which
adopts the GR--Athena++ form, Daszuta \emph{et al.} 2021, Eqs.~8--13, and the
Bona--Mass\'o\,$+\,\Gamma$--driver gauge, their Eq.~22). The evolved fields are
\begin{equation}
  \bm{u}=\bigl(\chi,\ \gt_{ij},\ \hat K,\ \At_{ij},\ \Gamt^{i},\ \Theta,\
  \alpha,\ \beta^{i},\ B^{i}\bigr),
  \label{eq:fields}
\end{equation}
with $\chi=(\det\gamma)^{\,p/12}$, $p\equiv\code{chi\_psi\_power}$
(default $p=-4$, so $\chi=\psi^{-2}$), $\gt_{ij}$ the unit--determinant
conformal metric, $\hat K=K-2\Theta$, and $\At_{ij}$ the conformal trace-free
extrinsic curvature. Schematically the evolution system is
\begin{equation}
  \partial_t \bm{u} = \Fop(\bm{u}),
  \label{eq:dtuF}
\end{equation}
which is LSK22 Eq.~(18). The concrete components, in the form actually coded
(see \S\ref{sec:standard}), are
\begin{align}
  \dtb\chi   &= \tfrac{p}{6}\,\chi\,(\alpha \hat K - \partial_k\beta^k),
    \label{eq:chi}\\
  \dtb\gt_{ij} &= -2\alpha\At_{ij}, \label{eq:g}\\
  \dtb\hat K &= -D^iD_i\alpha + \alpha\!\left(\At_{ij}\At^{ij}+\third\hat K^2\right)
        + \alpha\kappa_1(1-\kappa_2)\Theta + 4\pi\alpha(S+E),
    \label{eq:Khat}\\
  \dtb\At_{ij} &= \chi^{-4/p}\!\left[-D_iD_j\alpha+\alpha(R_{ij}+R^\phi_{ij})\right]^{\mathrm{TF}}
        + \alpha(\hat K\At_{ij}-2\At_{ik}\At^k{}_j)
        - 8\pi\alpha\!\left[\chi^{-4/p}S_{ij}\right]^{\mathrm{TF}},
    \label{eq:A}\\
  \dtb\Theta &= \alpha\!\left[\half\hat{\mathcal H} - (2+\kappa_2)\kappa_1\Theta\right]
        - 8\pi\alpha E,
    \label{eq:Theta}\\
  \dtb\Gamt^i &= 2\alpha\,\mathcal{D}^i
        - 2\At^{ij}\partial_j\alpha
        - 2\alpha\kappa_1(\Gamt^i-\Gamma^i_{(\gt)})
        - 16\pi\alpha\,\gt^{ij}S_j,
    \label{eq:Gam}
\end{align}
where $[\cdot]^{\mathrm{TF}}$ is the trace-free part, $\hat{\mathcal H}=R+\tfrac23\hat
K^2-\At_{ij}\At^{ij}$, $\mathcal D^i$ collects the $\At$/$\hat K$/$\Theta$
divergence terms, and the matter sources $(E,S_i,S_{ij},S=\gamma^{ij}S_{ij})$
come from the stress tensor. The gauge sector is
\begin{align}
  \partial_t\alpha &= \code{lapse\_advect}\,\beta^k\partial_k\alpha
       - f\,\alpha\,\hat K, \qquad
       f=\code{lapse\_oplog}\cdot\code{lapse\_harmonicf}+\code{lapse\_harmonic}\,\alpha,
    \label{eq:lapse}\\
  \partial_t\beta^i &= \code{shift\_Gamma}\,\Gamt^i
       + \code{shift\_advect}\,\beta^k\partial_k\beta^i
       - \code{shift\_eta}\,\beta^i + (\dots),
    \label{eq:shift}
\end{align}
which with the default parameters ($f=2$, $\code{shift\_Gamma}=1$,
$\code{shift\_eta}=2$) is the standard $1+\log$ slicing and (first--order)
$\Gamma$--driver of moving--puncture evolutions. Equations
\eqref{eq:lapse}--\eqref{eq:shift} are the analogue of LSK22 Eqs.~(8)--(10)
and Z24's Eq.~22 reference.

% ===========================================================================
\section{The continuum background--subtraction formulation}
% ===========================================================================

Following LSK22 (Eqs.~11--19) every evolved field is split into a
\emph{static analytic background} $\bg{\bm u}$ and a \emph{residual}
$\res{\bm u}$,
\begin{equation}
  \bm u = \bg{\bm u} + \res{\bm u}, \qquad \partial_t \bg{\bm u}=0,
  \label{eq:split}
\end{equation}
and the residual is evolved with the \emph{difference} of the RHS operator,
\begin{equation}
  \boxed{\ \partial_t \res{\bm u}
     = \Fop(\bg{\bm u}+\res{\bm u}) - \Fop(\bg{\bm u})\ }
  \label{eq:residualRHS}
\end{equation}
which is exactly LSK22 Eq.~(19). Two facts make this useful:
\begin{enumerate}[leftmargin=1.6em,itemsep=2pt]
\item \textbf{Continuum consistency.} If $\bg{\bm u}$ solves
$\Fop(\bg{\bm u})=0$ (a stationary solution of the \emph{evolution + gauge}
system), then \eqref{eq:residualRHS} reduces to $\partial_t\res{\bm
u}=\Fop(\bg{\bm u}+\res{\bm u})$, i.e.\ the residual carries the full nonlinear
dynamics with no spurious source.
\item \textbf{Discrete truncation cancellation (the key point).} The crucial
implementation requirement, emphasised in LSK22 immediately after their
Eq.~(19), is that \emph{the same discrete operator} $\Fnum$ be applied to both
arguments:
\begin{equation}
  \partial_t \res{\bm u}
     = \Fnum(\bg{\bm u}+\res{\bm u}) - \Fnum(\bg{\bm u}).
  \label{eq:residualRHSdiscrete}
\end{equation}
Because $\Fnum(\bg{\bm u})\neq 0$ numerically (it is the truncation error of
finite--differencing the smooth background), subtracting it removes the
$O(\res{\bm u}^0)$ background truncation error from the residual RHS. In
particular, if $\res{\bm u}\equiv 0$ then the two arguments are bitwise
identical and \eqref{eq:residualRHSdiscrete} is \emph{exactly} zero to
round--off: the static background truncation error never sources the residual.
\end{enumerate}

The remainder of this report verifies that \eqref{eq:residualRHSdiscrete} is
what the code computes, with one stencil, for \emph{every} field and every
auxiliary operation.

% ===========================================================================
\section{Code architecture and data flow}
% ===========================================================================

\subsection{State containers}
\file{z4c.hpp} (lines 80--133) declares four \code{Z4c\_vars} aliases over four
allocations, all of size \code{nz4c}:
\begin{center}
\begin{tabular}{lll}
\toprule
alias & allocation & meaning \\
\midrule
\code{z4c} & \code{u0} & evolved variable; in bg mode this holds the \emph{residual} $\res{\bm u}$\\
\code{bg}  & \code{u\_bg} & analytic background $\bg{\bm u}$ (recomputed each evaluation)\\
\code{full}& \code{u\_full} & reconstructed full state $\bg{\bm u}+\res{\bm u}$\\
\code{rhs} & \code{u\_rhs} & RHS storage $\partial_t\res{\bm u}$\\
\bottomrule
\end{tabular}
\end{center}
There is also a separate ADM background \code{adm\_bg} (\code{u\_adm\_bg}, 17
slots) used as the staging area for the analytic ADM background. The control
flags are (\file{z4c.cpp} 183--192):
\code{use\_analytic\_background}, \code{evolve\_gauge\_residual},
\code{evolve\_lapse\_residual}, \code{evolve\_shift\_residual},
\code{preserve\_lapse\_residual}. The function pointers
\code{SetADMBackground} and \code{SetZ4cBackground} are installed by the pgen.

\subsection{Transforms (\file{z4c.cpp})}
\begin{itemize}[leftmargin=1.4em,itemsep=2pt]
\item \code{ReconstructFullState()} (373--411): $\code{full}=\code{bg}+\code{z4c}$
field by field over the \emph{full} index range including ghosts
(\code{isg..ieg}). Lapse and shift residuals are added only if their evolve
(or, for the lapse, \code{preserve}) flag is set; otherwise the full gauge
equals the background gauge.
\item \code{RecastResidualState()} (413--452): the inverse,
$\code{z4c}=\code{full}-\code{bg}$, with the same gating.
\item \code{PrescribeGaugeResidual()} (454--484): zeroes $\res\alpha$,
$\res\beta^i$, $\res B^i$ when their residuals are not being evolved.
\item \code{UpdateBackgroundState()} (\file{z4c\_adm.cpp} 190--213): calls
\code{SetADMBackground} (fills \code{adm\_bg}), then either the closed--form
\code{SetZ4cBackground} (preferred) or the FD \code{ADMToZ4cImpl} fallback to
fill \code{bg}, and finishes with \code{EnforceAlgConstrOn(bg)}.
\end{itemize}

\subsection{Task ordering (\file{z4c\_tasks.cpp})}
The per--stage pipeline is
\code{CopyU}\,$\to$\,\code{CalcRHS}\,$\to$\,\code{Z4cBoundaryRHS}\,$\to$\,
\code{ExpRKUpdate}\,$\to$\,\code{RestrictU}\,$\to$\,\code{SendU/RecvU}\,$\to$\,
\code{ApplyPhysicalBCs}\,$\to$\,\code{Prolongate}\,$\to$\,\code{EnforceAlgConstr}\,$\to$\,
\code{ConvertZ4cToADM}. Two consequences matter for this audit:
\begin{itemize}[leftmargin=1.4em,itemsep=1pt]
\item All communication/refinement operators act on \code{u0}, i.e.\ on the
\emph{residual}. This is the desired behaviour (\S\ref{sec:amr}).
\item \code{EnforceAlgConstr} (194--204) reconstructs the full state,
projects the algebraic constraints \emph{on the full state}, and recasts the
residual (\S\ref{sec:alg}). \code{ConvertZ4cToADM} sets the ADM variables from
the full state, so matter, constraints, Weyl, and trackers all see the
physical metric.
\end{itemize}

% ===========================================================================
\section{The standard (non--background) RHS is unchanged}
\label{sec:standard}
% ===========================================================================

The main--branch monolithic kernel was refactored into two device functions in
\file{z4c\_calcrhs.cpp}:
\code{ComputeGeometryData} (127--398), which builds all derivatives, the
inverse metric, Christoffels, $R_{ij}$, $R^\phi_{ij}$, $D_iD_j\alpha$,
$\At$--contractions, $\mathcal D^i$, and the Lie/advection terms from a
templated \code{State}; and \code{BuildStandardPointwiseRHS} (400--500), which
assembles \eqref{eq:chi}--\eqref{eq:shift}. I verified line by line that
\code{BuildStandardPointwiseRHS} reproduces the \file{origin/main} kernel:
e.g.\ $\hat K$ (calcrhs 409--416 vs.\ main 482--489), $\chi$ (418--419 vs.\
490--491), $\Theta$ (421--427 vs.\ 492--499), $\Gamt^i$ (429--440 vs.\
501--513), $\gt_{ij}/\At_{ij}$ (442--459 vs.\ 516--532), lapse (461--464 vs.\
534--537), shift (483--499 vs.\ 540--551). The non--background path
(calcrhs 674--693) simply copies \code{BuildStandardPointwiseRHS(z4c,\dots)}
into \code{rhs}, so with \code{use\_analytic\_background=false} the branch is
numerically identical to \file{main} except for two \emph{additive} new gauge
options that default off (\code{slow\_start\_lapse}, \code{telegraph\_lapse},
\code{sss\_damping}) and the optional Gaussian rolling of $\kappa_1$
(\code{roll\_kappa}, calcrhs 540--549). \textbf{Conclusion:} the refactor is
behaviour--preserving for the full--field code.

% ===========================================================================
\section{Verification of the residual RHS, term by term}
\label{sec:rhs}
% ===========================================================================

With \code{use\_analytic\_background=true}, \code{CalcRHS} first refreshes the
background and reconstructs the full state (calcrhs 528--532),
\begin{lstlisting}
if (use_analytic_background) {
  pz4c->PrescribeGaugeResidual();
  pz4c->UpdateBackgroundState(time);
  pz4c->ReconstructFullState();
}
\end{lstlisting}
then, inside the point loop, computes \emph{two} geometry structures with the
\emph{same} \code{ComputeGeometryData<NGHOST>} stencil---one from
\code{full}, one from \code{bg} (calcrhs 556--567):
\begin{lstlisting}
ComputeGeometryData<NGHOST>(full, ..., geo_full);   // matter ON
ComputeGeometryData<NGHOST>(bg,   ..., geo_bg);      // matter OFF
\end{lstlisting}
The residual RHS is then assembled as the explicit difference. The mapping to
\eqref{eq:residualRHSdiscrete} is exact for every field:

{\small
\begin{longtable}{p{0.10\linewidth} p{0.50\linewidth} p{0.30\linewidth}}
\toprule
field & coded residual RHS $=\Fnum(\text{full})-\Fnum(\text{bg})$ & lines \\
\midrule
$\hat K$ &
$-(D^2\alpha|_f-D^2\alpha|_b)
 +\bigl(\alpha_f(\At\At+\third\hat K^2)_f-\alpha_b(\cdots)_b\bigr)
 +(\Lie\hat K|_f-\Lie\hat K|_b)
 +\kappa_1(1{-}\kappa_2)(\alpha_f\Theta_f-\alpha_b\Theta_b)$;
 matter $+4\pi\alpha_f(S_f+E)$ added to full only &
571--582 \\
\addlinespace
$\chi$ &
$(\Lie\chi|_f-\Lie\chi|_b)-\tfrac{p}{6}(\chi_f\alpha_f\hat K_f-\chi_b\alpha_b\hat K_b)$ &
584--588 \\
\addlinespace
$\Theta$ &
$(\Lie\Theta|_f-\Lie\Theta|_b)+\half(\alpha_f\hat{\mathcal H}_f-\alpha_b\hat{\mathcal H}_b)
 -(2{+}\kappa_2)\kappa_1(\alpha_f\Theta_f-\alpha_b\Theta_b)$;
 matter $-8\pi\alpha_f E$ to full only; $\times\,\code{use\_z4c}$ &
590--599 \\
\addlinespace
$\Gamt^i$ &
$2(\alpha_f\mathcal D^i_f-\alpha_b\mathcal D^i_b)+(\Lie\Gamt^i|_f-\Lie\Gamt^i|_b)
 -2\kappa_1(\alpha_f(\Gamt^i{-}\Gamma^i)_f-\alpha_b(\Gamt^i{-}\Gamma^i)_b)
 -2(\At^{ij}\partial_j\alpha|_f-|_b)$; matter to full only &
617--634 \\
\addlinespace
$\gt_{ij}$ &
$-2(\alpha_f\At_{ij,f}-\alpha_b\At_{ij,b})+(\Lie\gt_{ij}|_f-\Lie\gt_{ij}|_b)$ &
645--648 \\
\addlinespace
$\At_{ij}$ &
$\chi_f^{-4/p}[-D_iD_j\alpha+\alpha R]_f-\chi_b^{-4/p}[\cdots]_b
 -\third(\gt_{ij,f}[-D^2\alpha+\alpha R]_f-\gt_{ij,b}[\cdots]_b)
 +(\alpha_f(\hat K\At-2\At\At)_f-\alpha_b(\cdots)_b)+(\Lie\At|_f-\Lie\At|_b)$;
 matter to full only &
650--673 \\
\addlinespace
$\alpha$ &
$\code{rhs\_full\_gauge.alpha}-\code{rhs\_bg\_gauge.alpha}$ (else $0$) &
611--615 \\
\addlinespace
$\beta^i,B^i$ &
$\code{rhs\_full\_gauge.beta}-\code{rhs\_bg\_gauge.beta}$ (else $0$) &
635--641 \\
\bottomrule
\end{longtable}
}

\noindent For the gauge fields the difference is formed by calling
\code{BuildStandardPointwiseRHS} twice---once on \code{(full,geo\_full)} and
once on \code{(bg,geo\_bg)}---and subtracting (calcrhs 601--610). Thus the
gauge sector also realises \eqref{eq:residualRHSdiscrete} exactly.

\subsection{The zero--residual property holds}
Set $\res{\bm u}=0$. Then \code{ReconstructFullState} gives
$\code{full}\equiv\code{bg}$ bitwise. \code{ComputeGeometryData} is a pure
function of its \code{State}, so $\code{geo\_full}\equiv\code{geo\_bg}$ bitwise,
including every floored quantity (\code{chi\_guarded} uses the \emph{same}
\code{chi\_div\_floor} for both, calcrhs 276; \code{oopsi4} likewise). Every
table entry above is then a difference of identical numbers and evaluates to
$0$ to round--off. \textbf{This is the property the user asked about: no
truncation error of the (static) background propagates into the residual
evolution.} It holds because

\begin{enumerate}[leftmargin=1.6em,itemsep=1pt]
\item full and background geometry use one and the same discrete operator
(\code{ComputeGeometryData<NGHOST>}); there is no place where the background
is differenced at a different order or replaced by an analytic time
derivative;
\item the algebraic projection, dissipation, Sommerfeld BC, and AMR transfer
all act on the residual and inherit the same property (\S\S\ref{sec:alg}--\ref{sec:amr});
\item the matter source is added to the full evaluation only, which is correct
\emph{provided the background is vacuum} (\S\ref{sec:matter}).
\end{enumerate}

Away from $\res{\bm u}=0$, expanding \eqref{eq:residualRHSdiscrete} gives
$\partial_t\res{\bm u} = D\Fnum(\bg{\bm u})\!\cdot\!\res{\bm u} +
O(\res{\bm u}^2)$: the background truncation error enters only through the
\emph{coefficients} of the linearised operator (at $O(\res{\bm u}^1)$ and
higher), never as an additive $O(\res{\bm u}^0)$ source. Because Z4c is
quasilinear, the principal symbol of $D\Fnum(\bg{\bm u})$ equals that of the
full Z4c system linearised about $\bg{\bm u}$; strong hyperbolicity is
therefore inherited from the standard scheme (see \S\ref{sec:gauge} for the
gauge sector).

% ===========================================================================
\section{Background construction and its accuracy}
\label{sec:bgconstruct}
% ===========================================================================

\subsection{Closed--form Z4c background (preferred path)}
For the TOV/Kerr--Schild problem the pgen installs both
\code{SetADMBackground = SetADMBackgroundKerrSchild} and (when
\code{use\_direct\_z4c\_background=true}, the default)
\code{SetZ4cBackground = SetZ4cBackgroundKerrSchild}
(\file{z4c\_tov\_ks.cpp} 1688--1690). The latter (882--1038) fills
\emph{every} Z4c background field analytically, including the conformal
connection $\bg{\Gamt}^i$, which it computes from the closed--form metric
derivatives \code{ComputeMetricDerivatives}
\begin{equation}
  \bg{\Gamt}^i = -\partial_j\bigl(\psi^{4}\,\gt^{ij}\bigr)\Big|_{\text{analytic}}
  = -\bigl(\partial_j\psi^4\,\gt^{ij} + \psi^4\,\partial_j\gt^{ij}\bigr),
  \label{eq:Gambg}
\end{equation}
(z4c\_tov\_ks 1000--1036). This is important: it means $\bg{\Gamt}^i$ carries
\emph{no} grid--dependent truncation error and is identical on coarse and fine
AMR levels, which removes a resolution--dependent inconsistency that would
otherwise appear at refinement boundaries (\S\ref{sec:amr}).

\subsection{FD fallback}
If \code{use\_direct\_z4c\_background=false}, \code{UpdateBackgroundState}
falls back to \code{ADMToZ4cImpl} (\file{z4c\_adm.cpp} 107--116), which
computes $\bg{\Gamt}^i=-\partial_j\bg{\gt}^{ij}$ by finite differences. The
zero--residual cancellation of \S\ref{sec:rhs} \emph{still holds} (both
\code{geo\_full} and \code{geo\_bg} use the same stencil), but now
$\bg{\Gamt}^i$ itself is resolution dependent, so the background sampled on a
fine grid differs from the prolongation of the coarse--grid background at
$O(h^{\,n})$. \textbf{Recommendation:} keep
\code{use\_direct\_z4c\_background=true} for AMR runs (\S\ref{sec:robust}, R3).

\subsection{Algebraic consistency of the background}
\code{SetZ4cBackgroundKerrSchild} builds $\bg{\gt}_{ij}=(\det g)^{-1/3}g_{ij}$
and $\bg\chi=(\det g)^{p/12}$ so that $\det\bg{\gt}=1$ analytically, and
\code{UpdateBackgroundState} additionally applies
\code{EnforceAlgConstrOn(bg)} (\file{z4c\_adm.cpp} 198/212). Thus the
background satisfies the same algebraic constraints as the evolved state,
which is required for the projection in \S\ref{sec:alg} to leave a zero
residual invariant.

% ===========================================================================
\section{Gauge sector: is it a background--adapted $1+\log$?}
\label{sec:gauge}
% ===========================================================================

\subsection{What the code evolves}
With \code{evolve\_lapse\_residual = evolve\_shift\_residual = true} (the
production setting, see the sample deck, lines 113--123:
\code{lapse\_oplog=2}, \code{lapse\_harmonic=0}, \code{shift\_Gamma=1},
\code{shift\_advect=1}, \code{shift\_eta=2}), the residual lapse and shift obey
\begin{align}
  \partial_t\res\alpha
    &= \Bigl[\code{adv}\,\beta_f^k\partial_k\alpha_f - f_f\,\alpha_f\hat K_f\Bigr]
     - \Bigl[\code{adv}\,\beta_b^k\partial_k\alpha_b - f_b\,\alpha_b\hat K_b\Bigr],
    \label{eq:lapseres}\\
  \partial_t\res\beta^i
    &= \Bigl[\code{sG}\,\Gamt_f^i+\code{adv}\,\beta_f^k\partial_k\beta_f^i-\eta\beta_f^i\Bigr]
     - \Bigl[\code{sG}\,\Gamt_b^i+\code{adv}\,\beta_b^k\partial_k\beta_b^i-\eta\beta_b^i\Bigr].
    \label{eq:shiftres}
\end{align}
These are precisely \eqref{eq:residualRHSdiscrete} applied to
\eqref{eq:lapse}--\eqref{eq:shift}. By construction $\res\alpha=\res\beta^i=0$
is a fixed point: the background slicing/shift is stationary under the residual
gauge flow. This \emph{is} a background--adapted $1+\log$/$\Gamma$--driver,
matching the philosophy of LSK22 (their gauge Eqs.~8--10 enter the split
exactly as the geometric sector does).

\subsection{Principal symbol is preserved}
Because $\bg{\bm u}$ is time independent, subtracting the bracketed background
terms in \eqref{eq:lapseres}--\eqref{eq:shiftres} removes only
spatially--varying source/coefficient pieces; the highest derivative acting on
the \emph{residual} is unchanged from standard $1+\log$/$\Gamma$--driver,
namely the advection $\code{adv}\,\beta_f^k\partial_k(\cdot)$ and the
$\alpha\hat K$ / $\Gamt^i$ couplings. Hence the residual gauge subsystem has
the same principal symbol as the moving--puncture gauge and remains strongly
hyperbolic. The lapse advection in \eqref{eq:lapseres} uses $\beta_f$ for the
full term and $\beta_b$ for the background term; when the shift residual is off
($\beta_f=\beta_b=\bg\beta$) it reduces to advection of $\res\alpha$ by the
background shift, which is the consistent linearisation.

\subsection{The background is a fixed point; the only real difference from LSK22
is interior regularity}
\label{sec:gaugecaveat}
It is worth stating clearly what is, and what is \emph{not}, a concern here,
because it is easy to overstate the difference from LSK22.

\paragraph{The background \emph{is} an exact fixed point (no caveat).}
Equations \eqref{eq:lapseres}--\eqref{eq:shiftres} are, by construction,
$\partial_t\res\alpha=F_\alpha(\bg{\bm u}+\res{\bm u})-F_\alpha(\bg{\bm u})$ and
likewise for the shift, so $\res\alpha=\res\beta^i=0$ is an \emph{exact} fixed
point of the residual gauge flow at any resolution. This is the intended
behaviour and is correct.

\paragraph{The fixed point is stable; the moving--puncture damping still acts.}
Linearising \eqref{eq:lapseres} about $\res{\bm u}=0$ gives, schematically,
\begin{equation}
  \partial_t\res\alpha = \code{adv}\,\bg\beta^k\partial_k\res\alpha
     \;-\; f_0\,\bg\alpha\,\res{\hat K}
     \;-\; \underbrace{f_0\,\bg{\hat K}\,\res\alpha}_{\text{reaction}}
     \;+\;(\text{lower order}),
  \label{eq:lapselin}
\end{equation}
and the shift retains its $-\eta\,\res\beta^i$ damping and the constraint
damping $\kappa_1$ from the geometric sector. The principal part is identical to
standard $1+\log$/$\Gamma$--driver (\S\ref{sec:gauge}), so well--posedness and
strong hyperbolicity are unchanged; the \emph{new} terms (those proportional to
$\bg{\hat K}$, $\partial\bg\beta$, $\bg{\Gamt}^i$) are subprincipal. For
(non--extremal) Kerr--Schild $\bg{\hat K}>0$ outside the horizon, so the
reaction term in \eqref{eq:lapselin} is \emph{damping}, not growing. There is
thus no evidence that using Kerr--Schild instead of a trumpet introduces a new
gauge instability: the same damping channels ($\eta$, $\kappa_1$, the
$1+\log$ reaction) act on the residual.

\paragraph{What genuinely differs from LSK22.} LSK22 deliberately use a $K_0=0$
\emph{trumpet} puncture (their Eqs.~20--24), which is a continuum fixed point of
the puncture gauge, $\Fop(\bg{\bm u})=0$. Two consequences follow that are
worth keeping in mind, but neither is a stability problem:
\begin{enumerate}[leftmargin=1.6em,itemsep=2pt]
\item \emph{Clean vs.\ modified subtraction.} For the geometric/constraint
sector Kerr--Schild is a vacuum solution, so $\Fop_{\rm cont}(\bg{\bm u})=0$ and
the discrete subtraction removes \emph{only} truncation error---exactly as in
LSK22. For the \emph{gauge} sector $\Fop_{\rm cont}(\bg{\bm u})\neq0$, so the
subtraction additionally removes a nonzero continuum term, i.e.\ it turns
$1+\log$ into the background--adapted condition
$\partial_t\alpha=[1{+}\log](\alpha)-[1{+}\log](\bg\alpha)$ whose stationary
slice is Kerr--Schild rather than the trumpet. This is the deliberate
``background--adapted'' design, not a defect; the only practical implication is
that the reconstructed full field is \emph{not} a literal $1+\log$ evolution, so
gauge--dependent cross--comparisons with a standard moving--puncture run are not
apples--to--apples.
\item \emph{Interior regularity (the one tangible trade--off).} The trumpet
slice has a finite throat ($R=3M/2$ for Schwarzschild) and is regular over the
whole hypersurface, which is what makes moving--puncture evolutions
excision--free. Kerr--Schild is horizon penetrating and the slice reaches the
physical singularity at $r=0$, where the analytic background and its derivatives
diverge; the linearised operator $D\Fnum(\bg{\bm u})$ then has divergent
coefficients in the interior. The code handles this by (i) freezing the
background inside \code{excision\_freeze\_radius} (\file{z4c\_tov\_ks.cpp}
942--947) and (ii) ramping the residual RHS/state to zero inside
$\sim0.85\,r_{\rm hor}$ via \code{ApplyInnerExcision} (362--395). So an interior
excision/damping layer is \emph{required} for a Kerr--Schild background, whereas
a trumpet background would not need one. This is the substantive, and only,
robustness item; see R1 in \S\ref{sec:robust}.
\end{enumerate}

% ===========================================================================
\section{Auxiliary operations}
% ===========================================================================

\subsection{Algebraic constraints projected on the full state}
\label{sec:alg}
\code{EnforceAlgConstr} (\file{z4c\_tasks.cpp} 194--204) does
\code{ReconstructFullState}\,$\to$\,\code{EnforceAlgConstrOn(full)}\,$\to$\,
\code{RecastResidualState}. The unit--determinant and trace--free projections
(\file{z4c.cpp} 331--357) are \emph{nonlinear}, so they must be applied to the
physical conformal metric $\gt_{ij}=\bg{\gt}_{ij}+\res{\gt}_{ij}$, not to the
residual---which the code does correctly. Since the background is itself
algebraically constrained (\S\ref{sec:bgconstruct}), $\res{\bm u}=0\Rightarrow
\code{full}=\code{bg}$ already satisfies the constraints, the projection is the
identity, and the recast leaves $\res{\bm u}=0$. \textbf{Correct.}

\subsection{Kreiss--Oliger dissipation acts on the residual}
\label{sec:diss}
The dissipation loop (\file{z4c\_calcrhs.cpp} 702--709) adds
$\code{diss}\cdot\mathrm{Diss}(\code{u0})$ to \code{u\_rhs} for all
\code{nz4c} fields, where \code{u0} is the residual. This is the right choice:
dissipating the \emph{full} field would inject the smooth background's
high--wavenumber truncation content into the evolution, whereas dissipating the
residual vanishes when $\res{\bm u}=0$ and only smooths genuine residual
features. \textbf{Correct and beneficial.} One minor inconsistency: if
\code{preserve\_lapse\_residual=true} while \code{evolve\_lapse\_residual=false},
then $\partial_t\res\alpha$ from \code{CalcRHS} is $0$ (calcrhs 614) but the
dissipation still adds $\code{diss}\cdot\mathrm{Diss}(\res\alpha)$, so a
``preserved'' lapse residual slowly diffuses rather than being held fixed
(R4).

\subsection{Sommerfeld boundary on the residual}
\label{sec:sbc}
\code{Z4cBoundaryRHS}/\code{Z4cSommerfeld} (\file{z4c\_Sbc.cpp}) overwrite the
outflow--boundary RHS of $\hat K,\Theta,\Gamt^i,\At_{ij}$ with the standard
radiative form $\partial_t u = -v\,(s^k\partial_k u + u/r)$ using the
\code{z4c} alias (the \emph{residual}). For the residual this is the correct
asymptotic condition: the residual $\to 0$ at large $r$ where the background is
asymptotically flat, so $u_{\rm asympt}=0$ and the constant--coefficient
speeds ($v=1$, and $v=\sqrt2$ for $\hat K$) apply. When $\res{\bm u}=0$ the
Sommerfeld RHS is $0$, consistent with the zero--residual property.
\textbf{Correct.} Caveat: the speeds assume near--flatness at the boundary; for
boxes whose outer faces sit at moderate $r$ (e.g.\ the $-16$ face in the sample
deck) there is a small mismatch, common to all Z4c outflow BCs and not
specific to background subtraction (R5).

\subsection{AMR refinement and transfer on the residual}
\label{sec:amr}
Communication (\code{SendU/RecvU}), restriction (\code{RestrictU}), and
prolongation (\code{Prolongate}) all operate on \code{u0} = residual. Because
the residual is the smooth physical perturbation (the steep Kerr--Schild
gradients live in the analytic background, which is regenerated on every
level), interpolation/restriction error is dramatically smaller than it would
be for the full field---this is one of the main benefits of the method and is
consistent with Z24's higher--order CC prolongation/restriction. The
$\chi$/$d\chi$ AMR criteria were correctly updated to refine on the
\emph{reconstructed full} conformal factor (\file{z4c\_amr.cpp}: \code{Refine}
reconstructs the full state when the method is \code{Chi}/\code{dChi};
\code{RefineChiMin}/\code{RefineDchiMax} read \code{u\_full}), so the BH is
tracked by the full $\chi$ rather than by the residual. \textbf{Correct.}

\subsection{Matter coupling and the vacuum--background assumption}
\label{sec:matter}
The matter sources $4\pi\alpha(S+E)$, $-8\pi\alpha E$, $-16\pi\alpha\gt^{ij}S_j$,
$-8\pi\alpha[\chi^{-4/p}S_{ij}]^{\rm TF}$ are added to the \emph{full}
evaluation only (calcrhs 579--581, 596--598, 630--633, 667--672). This is
correct \emph{iff the analytic background carries no matter}, which is the case
here: the background is the vacuum (Kerr--Schild or Minkowski) black--hole
spacetime, and the TOV star is the matter that sources the residual
(\file{z4c\_tov\_ks.cpp}; the BH background is built by
\code{ComputeADMDecomposition}/\code{FillFlatADM} with $K^{\rm
matter}_{ij}=0$ and no stress tensor). \textbf{Correct for the intended use.}
\emph{Hazard:} if a future background ever includes its own matter (e.g.\ using
an equilibrium TOV \emph{as the background}), the source would be double
counted and $\res{\bm u}=0$ would no longer be a fixed point; one would then
have to subtract the background matter source as well. See R2.

\subsection{Initial data and restart}
\code{ConvertADMToResidualOnBackground} (\file{z4c\_tov\_ks.cpp} 1307--1349)
constructs the full Z4c state by ADM$\to$Z4c, builds a flat reference, forms
the difference, then the setup (1687--1697) regenerates the analytic
Kerr--Schild background, reconstructs full $=$ bg $+$ residual, projects the
algebraic constraints on the full state, and recasts the residual---so the
stored residual is finally defined with respect to the same analytic
background used during evolution. The restart branch (1886--1896) repeats the
background regeneration / reconstruction / recast, so a restart is consistent
with a continued run. \textbf{Correct.}

% ===========================================================================
\section{Robustness assessment (ranked)}
\label{sec:robust}
% ===========================================================================

\newcommand{\sev}[1]{\textbf{[#1]}}

\begin{description}[leftmargin=2.2em,style=unboxed,itemsep=4pt]

\item[R1 \sev{Design trade--off / Low--Medium}] \emph{Interior regularity:
Kerr--Schild reaches the singularity, so an interior excision/damping layer is
required.} \emph{This is \textbf{not} a statement that the background fails to be
a fixed point---it is one, exactly and stably} (\S\ref{sec:gaugecaveat}). The
real difference from LSK22 is that their $K_0=0$ trumpet has a finite throat and
is regular over the whole slice (hence excision--free), whereas the
Kerr--Schild background is horizon penetrating and diverges at $r=0$, so the
linearised operator has divergent interior coefficients. The code copes by
freezing the background inside \code{excision\_freeze\_radius} and ramping the
residual RHS/state to zero inside $\sim0.85\,r_{\rm hor}$
(\code{ApplyInnerExcision}, \file{z4c\_tov\_ks.cpp} 362--395, 942--947); this
layer is load--bearing for Kerr--Schild and its parameters should be treated as
part of the scheme. A secondary, benign point: because the gauge subtraction
removes a nonzero \emph{continuum} term, the reconstructed full field obeys a
background--adapted $1+\log$ (stationary at Kerr--Schild) rather than literal
$1+\log$, so gauge--dependent comparisons to standard moving--puncture runs are
not apples--to--apples. \textbf{Mitigations:} (i) if a fully interior--regular,
excision--free setup is wanted, use a $K_0=0$ trumpet background as in LSK22;
(ii) verify robustness of the excision parameters in a convergence test and
monitor the already--exposed gauge residuals (\code{alpha-res}, \code{beta-res},
\code{B-res}, \code{Gam-res}; \file{z4c\_tov\_ks.cpp} 421--424).

\item[R2 \sev{By design / Cosmetic}] \emph{Matter source assumes a vacuum
background.} In the intended use the matter (the TOV star, and any accreting
gas) is \emph{always} part of the perturbation and the analytic background is
always a vacuum black hole, so the vacuum--background assumption holds by
construction and this item never fires. The only un--addressed case is a
\emph{deliberate} future change that puts matter into the background itself
(e.g.\ using an equilibrium TOV as the background), which would double--count
the source (\S\ref{sec:matter}); there is currently no guard against that
misuse. \textbf{Optional mitigation (defensive only):} add an assertion that the
background stress tensor vanishes, so the method cannot be silently misused.

\item[R3 \sev{Accuracy / Low--Medium}] \emph{FD background connection under
AMR.} Only one background field, the conformal connection $\bg{\Gamt}^i$,
is defined through a \emph{spatial derivative} ($\bg{\Gamt}^i=-\partial_j
\bg{\gt}^{ij}$); all the others are pointwise evaluations of the analytic
solution and are therefore resolution independent. With
\code{use\_direct\_z4c\_background=false}, $\bg{\Gamt}^i$ is finite--differenced
on the grid, so its value carries the stencil truncation $O(h^n)$ and
\emph{differs between coarse and fine levels at the same physical point}. Two
consequences follow: (a) the stored residual $\res{\Gamt}^i=\Gamt^i_{\rm
full}-\bg{\Gamt}^i$ is no longer a purely smooth physical field---it absorbs the
grid--dependent background truncation, partly defeating the smoothness that
makes residual prolongation/restriction accurate (\S\ref{sec:amr}); and (b) at a
refinement boundary the residual is defined relative to one level's
$\bg{\Gamt}^i$ but reconstructed against another's, so
$\code{full}=\res{}+\bg{}$ picks up an $O(h^n_{\rm coarse})-O(h^n_{\rm fine})$
mismatch exactly the kind of background truncation the method is meant to
cancel. With the default closed--form path \eqref{eq:Gambg}, $\bg{\Gamt}^i$ is
analytic and bit--for--bit identical on every level, so neither effect occurs.
\textbf{Mitigation:} keep \code{use\_direct\_z4c\_background=true} (the default)
for all multilevel runs; if the FD fallback is ever required, recompute
$\bg{\Gamt}^i$ from the same analytic derivatives on every level.

\item[R4 \sev{Bug / Low}] \emph{Dissipation of a ``preserved'' lapse residual.}
With \code{preserve\_lapse\_residual=true} and
\code{evolve\_lapse\_residual=false}, $\res\alpha$ has zero analytic RHS but
still receives K--O dissipation (\S\ref{sec:diss}), so it is not actually held
fixed. \textbf{Mitigation:} skip the dissipation contribution for non--evolved
gauge components, or zero \code{u\_rhs} for those slots after the dissipation
loop.

\item[R5 \sev{Pre--existing / Low}] \emph{Telegraph--lapse auxiliary
$B$ is gated on the shift flag.} In the residual path, $\res B^i$ (used only
for \code{telegraph\_lapse}) is updated inside the
\code{if (evolve\_shift\_residual)} block (calcrhs 635--641), so a
telegraph lapse with shift residual disabled would not evolve its heat--flux
variable. Niche (telegraph lapse is off by default), but the gating is on the
wrong flag. \textbf{Mitigation:} gate $\res B^i$ on the lapse flag (or on
\code{telegraph\_lapse}).

\item[R6 \sev{Latent / Low}] \emph{Sub--stage time for time--dependent
backgrounds.} \code{CalcRHS} evaluates the background at
\code{pmesh->time} (calcrhs 522), i.e.\ the step--start time, for every RK
substage. For the current \emph{static} Kerr--Schild background this is exact
(and the pgen forces the BH to the origin, 1785--1792). For any future
time--dependent background it would introduce an $O(\Delta t)$ phase error that
the residual must absorb, partially defeating the subtraction.
\textbf{Mitigation:} pass the substage time if/when a moving background is
enabled.

\item[R7 \sev{Performance / Cosmetic}] \emph{Redundant background work.}
\code{UpdateBackgroundState}\,+\,\code{ReconstructFullState} run in
\code{CalcRHS}, again in \code{EnforceAlgConstr}, and again in
\code{ConvertZ4cToADM}/\code{ADMConstraints\_}/AMR within a single stage; and
\code{BuildStandardPointwiseRHS} is evaluated in full for the background even
though only its gauge outputs are used (calcrhs 607--609). Harmless (the
background is idempotent) but wasteful. \textbf{Mitigation:} cache the
background per (time, stage); compute only the gauge components for the
background.

\end{description}

\subsection*{Things that are explicitly correct}
\begin{itemize}[leftmargin=1.4em,itemsep=1pt]
\item Residual RHS $=\Fnum(\text{full})-\Fnum(\text{bg})$ for \emph{every}
field, with one shared stencil $\Rightarrow$ exact zero--residual property and
no additive background--truncation source (\S\ref{sec:rhs}).
\item Floors (\code{chi\_div\_floor}, \code{chi\_min\_floor}) and \code{oopsi4}
are identical between full and background, so they do not break the
cancellation.
\item Algebraic constraints projected on the \emph{full} state, then recast
(\S\ref{sec:alg}).
\item K--O dissipation, Sommerfeld BC, and AMR transfer all act on the
residual (\S\S\ref{sec:diss}--\ref{sec:amr}).
\item Closed--form analytic background including $\bg{\Gamt}^i$, removing
resolution dependence of the background (\S\ref{sec:bgconstruct}).
\item Background--adapted gauge retains the $1+\log$/$\Gamma$--driver principal
symbol (strong hyperbolicity inherited) (\S\ref{sec:gauge}).
\item Non--background code path is numerically identical to \file{origin/main}
(\S\ref{sec:standard}).
\item Constraints/Weyl/ADM/matter all see the reconstructed full physical
metric (\S\ref{sec:standard} task ordering).
\end{itemize}

% ===========================================================================
\section{Recommendations}
% ===========================================================================
\begin{enumerate}[leftmargin=1.6em,itemsep=2pt]
\item \textbf{Decide the background philosophy (R1).} The Kerr--Schild
background is an exact, stable fixed point of the residual gauge, so this is not
a correctness issue. The only tangible trade--off is interior regularity: if an
excision--free, interior--regular setup is desired, adopt a $K_0=0$ trumpet
background (LSK22 Eqs.~20--25); otherwise keep Kerr--Schild but treat the inner
freeze/excision layer as a load--bearing part of the scheme and verify its
parameter robustness (and gauge--residual stationarity) in convergence tests.
\item \textbf{Add a vacuum--background assertion (R2)} (or generalise the
matter subtraction) so the method cannot be silently misused with a
matter--bearing background.
\item \textbf{Keep \code{use\_direct\_z4c\_background=true} for AMR (R3)}; if
the FD fallback must be used, recompute $\bg{\Gamt}^i$ consistently across
levels.
\item \textbf{Fix the gauge--gating issues (R4, R5)} so that non--evolved /
preserved gauge components are neither dissipated nor mis--gated.
\item \textbf{Numerical zero--residual test.} Run with
\code{zero\_tmunu=true} (already supported, 1710--1715) and a pure background:
verify $\|\partial_t\res{\bm u}\|$ stays at round--off and is independent of
resolution---this directly certifies \eqref{eq:residualRHSdiscrete} and the
absence of background--truncation leakage.
\item \textbf{Convergence of the residual} at fixed physical perturbation to
confirm the expected order and that the subtraction does not degrade it.
\end{enumerate}

% ===========================================================================
\section{Summary}
% ===========================================================================
The background--subtracted Z4c in \file{src/z4c} faithfully implements LSK22
Eq.~(19), $\partial_t\res{\bm u}=\Fnum(\bg{\bm u}+\res{\bm u})-\Fnum(\bg{\bm
u})$, on top of the Z24/GR--Athena++ Z4c system. The subtraction is performed
field by field with a single shared finite--difference operator, so the
zero--residual property holds to round--off and the static background's
truncation error does not source the residual; the algebraic projection,
dissipation, boundary, AMR, and matter coupling are all consistent with this.
The gauge is a genuine background--adapted $1+\log$/$\Gamma$--driver that keeps
the standard principal symbol, and the Kerr--Schild background is an exact,
stable fixed point of the residual gauge flow (the moving--puncture damping
channels still act on the residual). The only real design trade--off is that
Kerr--Schild is horizon penetrating and reaches the singularity, so an interior
freeze/excision layer is required where a $K_0=0$ trumpet would be regular
(R1)---a deliberate choice, not a stability defect. The remaining items
(R2--R7) are latent hazards, a minor dissipation/gating bug, and performance
redundancies.

\end{document}
